\chapter{Casos de uso - Tercera iteración}
\label{chap:anexo_casos_uso_it3}

Este anexo presenta la especificación detallada de los casos de uso desarrollados durante la tercera iteración del sistema, enfocados en la arquitectura API completa y funcionalidades avanzadas.

\section{CU-03 - Desarrollo de aplicaciones BCI personalizadas}
\label{sec:cu03_it3_desarrollo_bci}

\begin{table}[H]
\centering
\resizebox{\textwidth}{!}{
\rowcolors{2}{white}{udcgray!25}
\begin{tabular}{|p{3cm}|p{12cm}|}
\hline
\rowcolor{udcpink!25}
\textbf{Campo} & \textbf{Descripción} \\
\hline
\textbf{Identificador} & CU-03 \\
\hline
\textbf{Nombre} & Desarrollo de aplicaciones BCI personalizadas mediante acceso programático completo \\
\hline
\textbf{Actores} & Desarrollador de aplicaciones BCI (principal), Investigador con conocimientos técnicos (secundario) \\
\hline
\textbf{Descripción} & El sistema debe proporcionar APIs completas que permitan desarrollo de aplicaciones BCI personalizadas con acceso programático a todas las capacidades del sistema \\
\hline
\textbf{Precondiciones} & 
\begin{minipage}[t]{11cm}
\begin{itemize}
    \item Las APIs de acceso a datos están completamente funcionales
    \item Las APIs de control están documentadas y testadas
    \item El explorador web proporciona documentación interactiva
    \item Los endpoints de streaming en tiempo real están operativos
\end{itemize}
\end{minipage} \\
\hline
\textbf{Postcondiciones} & 
\begin{minipage}[t]{11cm}
\begin{itemize}
    \item La aplicación personalizada accede a datos en tiempo real
    \item El control remoto de dispositivos funciona correctamente
    \item La integración preserva funcionalidad del sistema base
    \item Los datos son compatibles con frameworks ML existentes
\end{itemize}
\end{minipage} \\
\hline
\end{tabular}
}
\caption{Especificación del caso de uso CU-03 - Tercera iteración}
\label{tab:cu03_it3_especificacion}
\end{table}

\subsection{Flujo principal}
\label{subsec:cu03_it3_flujo_principal}

\begin{table}[H]
\centering
\resizebox{\textwidth}{!}{
\rowcolors{2}{white}{udcgray!25}
\begin{tabular}{|c|p{13cm}|}
\hline
\rowcolor{udcpink!25}
\textbf{Paso} & \textbf{Acción} \\
\hline
1 & El desarrollador accede al explorador web para documentación API \\
\hline
2 & Explora endpoints disponibles mediante interfaz interactiva \\
\hline
3 & Testa funcionalidades específicas usando plantillas JSON proporcionadas \\
\hline
4 & Diseña aplicación personalizada utilizando APIs REST y WebSocket \\
\hline
5 & Implementa cliente que consume endpoints de acceso a datos \\
\hline
6 & Integra control remoto mediante APIs de interacción \\
\hline
7 & Establece streaming de datos en tiempo real via WebSocket \\
\hline
8 & Configura callbacks para eventos del sistema (conexión, errores, datos) \\
\hline
9 & Implementa lógica de procesamiento específica de la aplicación \\
\hline
10 & Valida funcionamiento con datos sintéticos y reales \\
\hline
11 & Integra con frameworks ML (TensorFlow, PyTorch) según necesidades \\
\hline
12 & Despliega aplicación con capacidades BCI completas \\
\hline
\end{tabular}
}
\caption{Flujo principal del caso de uso CU-03 - Tercera iteración}
\label{tab:cu03_it3_flujo_principal}
\end{table}

\subsection{Flujos alternativos}
\label{subsec:cu03_it3_flujos_alternativos}

\begin{table}[H]
\centering
\resizebox{\textwidth}{!}{
\rowcolors{2}{white}{udcgray!25}
\begin{tabular}{|p{3cm}|p{11cm}|}
\hline
\rowcolor{udcpink!25}
\textbf{Flujo} & \textbf{Descripción} \\
\hline
\textbf{FA-01: Desarrollo CLI personalizado} & 
\begin{minipage}[t]{10.5cm}
\begin{itemize}
    \item En paso 4, desarrollador crea herramientas CLI usando Typer
    \item Utiliza APIs internas para funcionalidades específicas
    \item Implementa comandos especializados para flujo de trabajo único
    \item Integra con pipelines de automatización existentes
\end{itemize}
\end{minipage} \\
\hline
\textbf{FA-02: Aplicación web personalizada} & 
\begin{minipage}[t]{10.5cm}
\begin{itemize}
    \item Desarrolla frontend web que consume APIs directamente
    \item Implementa visualizaciones específicas usando Chart.js/D3.js
    \item Crea dashboards personalizados para experimentos específicos
    \item Integra autenticación y gestión de usuarios si necesario
\end{itemize}
\end{minipage} \\
\hline
\textbf{FA-03: Integración framework específico} & 
\begin{minipage}[t]{10.5cm}
\begin{itemize}
    \item Desarrolla adaptadores para frameworks específicos (Unity, ROS)
    \item Implementa protocolos de comunicación nativos del framework
    \item Mantiene sincronización temporal específica del dominio
    \item Proporciona APIs nivel superior adaptadas al framework
\end{itemize}
\end{minipage} \\
\hline
\end{tabular}
}
\caption{Flujos alternativos del caso de uso CU-03 - Tercera iteración}
\label{tab:cu03_it3_flujos_alternativos}
\end{table}

\subsection{Flujos de excepción}
\label{subsec:cu03_it3_flujos_excepcion}

\begin{table}[H]
\centering
\resizebox{\textwidth}{!}{
\rowcolors{2}{white}{udcgray!25}
\begin{tabular}{|p{3cm}|p{11cm}|}
\hline
\rowcolor{udcpink!25}
\textbf{Excepción} & \textbf{Manejo} \\
\hline
\textbf{FE-01: Fallo WebSocket streaming} & 
\begin{minipage}[t]{10.5cm}
\begin{itemize}
    \item Interrupción en streaming de datos en tiempo real
    \item Sistema detecta desconexión y intenta reconexión automática
    \item Mantiene buffer de datos durante interrupción temporal
    \item Proporciona callback de error con detalles específicos
\end{itemize}
\end{minipage} \\
\hline
\textbf{FE-02: API endpoint no disponible} & 
\begin{minipage}[t]{10.5cm}
\begin{itemize}
    \item Servicio específico no responde o está temporalmente inactivo
    \item Sistema proporciona mensaje de error HTTP descriptivo
    \item Sugiere endpoints alternativos para funcionalidad similar
    \item Implementa retry automático con backoff exponencial
\end{itemize}
\end{minipage} \\
\hline
\textbf{FE-03: Datos malformados en API} & 
\begin{minipage}[t]{10.5cm}
\begin{itemize}
    \item Respuesta de API no cumple formato JSON esperado
    \item Sistema valida estructura antes de procesar datos
    \item Proporciona esquema esperado en mensaje de error
    \item Permite modo degradado con datos parciales válidos
\end{itemize}
\end{minipage} \\
\hline
\textbf{FE-04: Exceso límites rate limiting} & 
\begin{minipage}[t]{10.5cm}
\begin{itemize}
    \item Aplicación supera límites de frecuencia de peticiones
    \item Sistema implementa throttling automático inteligente
    \item Proporciona headers con información de límites restantes
    \item Sugiere optimizaciones para reducir frecuencia de peticiones
\end{itemize}
\end{minipage} \\
\hline
\end{tabular}
}
\caption{Flujos de excepción del caso de uso CU-03 - Tercera iteración}
\label{tab:cu03_it3_flujos_excepcion}
\end{table}

\section{CU-04 - Herramientas avanzadas para investigación neurocientífica}
\label{sec:cu04_it3_investigacion_avanzada}

\begin{table}[H]
\centering
\resizebox{\textwidth}{!}{
\rowcolors{2}{white}{udcgray!25}
\begin{tabular}{|p{3cm}|p{12cm}|}
\hline
\rowcolor{udcpink!25}
\textbf{Campo} & \textbf{Descripción} \\
\hline
\textbf{Identificador} & CU-04 \\
\hline
\textbf{Nombre} & Herramientas avanzadas para investigación neurocientífica con algoritmos especializados \\
\hline
\textbf{Actores} & Investigador neurocientífico (principal), Personal técnico de laboratorio (secundario) \\
\hline
\textbf{Descripción} & El sistema debe proporcionar algoritmos avanzados de procesamiento EEG como ICA, wavelets, e integración MNE-Python para análisis neurofisiológico completo \\
\hline
\textbf{Precondiciones} & 
\begin{minipage}[t]{11cm}
\begin{itemize}
    \item Los servicios avanzados están completamente implementados
    \item La integración MNE-Python está funcional
    \item Los algoritmos ICA y wavelet están optimizados
    \item Los visualizadores especializados están operativos
\end{itemize}
\end{minipage} \\
\hline
\textbf{Postcondiciones} & 
\begin{minipage}[t]{11cm}
\begin{itemize}
    \item Los análisis avanzados proporcionan resultados interpretables
    \item Las visualizaciones facilitan interpretación neurofisiológica
    \item Los resultados son exportables en formatos científicos
    \item La documentación facilita reproducibilidad de análisis
\end{itemize}
\end{minipage} \\
\hline
\end{tabular}
}
\caption{Especificación del caso de uso CU-04 - Tercera iteración}
\label{tab:cu04_it3_especificacion}
\end{table}

\subsection{Flujo principal}
\label{subsec:cu04_it3_flujo_principal}

\begin{table}[H]
\centering
\resizebox{\textwidth}{!}{
\rowcolors{2}{white}{udcgray!25}
\begin{tabular}{|c|p{13cm}|}
\hline
\rowcolor{udcpink!25}
\textbf{Paso} & \textbf{Acción} \\
\hline
1 & El investigador configura dispositivo EEG y verifica calidad de señal \\
\hline
2 & Accede a panel de procesamiento avanzado desde interfaz principal \\
\hline
3 & Configura filtros digitales especializados según protocolo experimental \\
\hline
4 & Selecciona canales específicos para análisis ICA (1-16 componentes) \\
\hline
5 & Ejecuta análisis ICA y visualiza matrices mixing/unmixing resultantes \\
\hline
6 & Analiza componentes individuales mediante weight distribution y PSD \\
\hline
7 & Configura transformadas wavelet (familia, niveles) según necesidades \\
\hline
8 & Aplica análisis tiempo-frecuencia mediante wavelets \\
\hline
9 & Configura conversión MNE-Python (modalidad: flexible/rendimiento/tiempo real) \\
\hline
10 & Genera visualizaciones científicas (topomaps, espectrogramas, PSD) \\
\hline
11 & Exporta resultados en formatos compatibles con publicación científica \\
\hline
12 & Documenta pipeline de análisis para reproducibilidad \\
\hline
\end{tabular}
}
\caption{Flujo principal del caso de uso CU-04 - Tercera iteración}
\label{tab:cu04_it3_flujo_principal}
\end{table}

\subsection{Flujos alternativos}
\label{subsec:cu04_it3_flujos_alternativos}

\begin{table}[H]
\centering
\resizebox{\textwidth}{!}{
\rowcolors{2}{white}{udcgray!25}
\begin{tabular}{|p{3cm}|p{11cm}|}
\hline
\rowcolor{udcpink!25}
\textbf{Flujo} & \textbf{Descripción} \\
\hline
\textbf{FA-01: Análisis ICA tiempo real} & 
\begin{minipage}[t]{10.5cm}
\begin{itemize}
    \item En paso 5, investigador activa modo tiempo real para ICA
    \item Sistema recalcula componentes automáticamente con nuevos datos
    \item Visualizaciones se actualizan dinámicamente con latencia aceptable
    \item Permite identificación inmediata de artefactos durante captura
\end{itemize}
\end{minipage} \\
\hline
\textbf{FA-02: Análisis multi-banda wavelet} & 
\begin{minipage}[t]{10.5cm}
\begin{itemize}
    \item En paso 8, se configura análisis simultáneo múltiples familias
    \item Sistema compara resultados Daubechies, Haar, Biorthogonal
    \item Genera visualización comparativa de descomposiciones
    \item Facilita selección óptima de wavelet para datos específicos
\end{itemize}
\end{minipage} \\
\hline
\textbf{FA-03: Pipeline automatizado MNE} & 
\begin{minipage}[t]{10.5cm}
\begin{itemize}
    \item En paso 9, se configura pipeline automatizado de análisis
    \item Sistema ejecuta secuencia predefinida: filtrado → ICA → análisis
    \item Genera reporte automático con todas las visualizaciones
    \item Exporta resultados en formato publication-ready
\end{itemize}
\end{minipage} \\
\hline
\end{tabular}
}
\caption{Flujos alternativos del caso de uso CU-04 - Tercera iteración}
\label{tab:cu04_it3_flujos_alternativos}
\end{table}

\subsection{Flujos de excepción}
\label{subsec:cu04_it3_flujos_excepcion}

\begin{table}[H]
\centering
\resizebox{\textwidth}{!}{
\rowcolors{2}{white}{udcgray!25}
\begin{tabular}{|p{3cm}|p{11cm}|}
\hline
\rowcolor{udcpink!25}
\textbf{Excepción} & \textbf{Manejo} \\
\hline
\textbf{FE-01: Fallo convergencia ICA} & 
\begin{minipage}[t]{10.5cm}
\begin{itemize}
    \item Algoritmo ICA no converge con parámetros especificados
    \item Sistema sugiere reducir número de componentes automáticamente
    \item Ofrece probar con datos pre-filtrados o mayor ventana temporal
    \item Proporciona diagnóstico específico sobre calidad de datos
\end{itemize}
\end{minipage} \\
\hline
\textbf{FE-02: Datos insuficientes análisis} & 
\begin{minipage}[t]{10.5cm}
\begin{itemize}
    \item Ventana temporal muy pequeña para análisis robusto
    \item Sistema calcula mínimo requerido según algoritmo específico
    \item Sugiere extender captura o usar datos históricos disponibles
    \item Permite análisis con advertencia sobre limitaciones estadísticas
\end{itemize}
\end{minipage} \\
\hline
\textbf{FE-03: Error integración MNE} & 
\begin{minipage}[t]{10.5cm}
\begin{itemize}
    \item Fallo en conversión o procesamiento MNE-Python
    \item Sistema identifica paso específico que causa error
    \item Proporciona información detallada sobre incompatibilidad
    \item Ofrece conversión alternativa o modo degradado sin MNE
\end{itemize}
\end{minipage} \\
\hline
\textbf{FE-04: Sobrecarga computacional} & 
\begin{minipage}[t]{10.5cm}
\begin{itemize}
    \item Análisis requiere más recursos que disponibles en sistema
    \item Sistema ofrece procesamiento batch en lugar de tiempo real
    \item Sugiere reducir complejidad o usar submuestreo de datos
    \item Proporciona estimación de tiempo requerido para análisis
\end{itemize}
\end{minipage} \\
\hline
\end{tabular}
}
\caption{Flujos de excepción del caso de uso CU-04 - Tercera iteración}
\label{tab:cu04_it3_flujos_excepcion}
\end{table}

\section{CU-06 - Integración programática con herramientas existentes}
\label{sec:cu06_it3_integracion_programatica}

\begin{table}[H]
\centering
\resizebox{\textwidth}{!}{
\rowcolors{2}{white}{udcgray!25}
\begin{tabular}{|p{3cm}|p{12cm}|}
\hline
\rowcolor{udcpink!25}
\textbf{Campo} & \textbf{Descripción} \\
\hline
\textbf{Identificador} & CU-06 \\
\hline
\textbf{Nombre} & Integración programática con herramientas de análisis existentes mediante APIs estándar \\
\hline
\textbf{Actores} & Integrador de sistemas (principal), Analista de datos (secundario) \\
\hline
\textbf{Descripción} & El sistema debe facilitar la integración programática con herramientas existentes mediante APIs documentadas y formatos de datos estándar \\
\hline
\textbf{Precondiciones} & 
\begin{minipage}[t]{11cm}
\begin{itemize}
    \item Las APIs están completamente documentadas
    \item Los formatos de exportación están estandarizados
    \item Los endpoints de streaming están estables
    \item La documentación OpenAPI está actualizada
\end{itemize}
\end{minipage} \\
\hline
\textbf{Postcondiciones} & 
\begin{minipage}[t]{11cm}
\begin{itemize}
    \item La integración funciona de forma transparente
    \item Los datos mantienen integridad durante transferencia
    \item Los pipelines automatizados operan sin intervención
    \item La documentación facilita mantenimiento futuro
\end{itemize}
\end{minipage} \\
\hline
\end{tabular}
}
\caption{Especificación del caso de uso CU-06 - Tercera iteración}
\label{tab:cu06_it3_especificacion}
\end{table}

\subsection{Flujo principal}
\label{subsec:cu06_it3_flujo_principal}

\begin{table}[H]
\centering
\resizebox{\textwidth}{!}{
\rowcolors{2}{white}{udcgray!25}
\begin{tabular}{|c|p{13cm}|}
\hline
\rowcolor{udcpink!25}
\textbf{Paso} & \textbf{Acción} \\
\hline
1 & El integrador analiza documentación API mediante explorador web \\
\hline
2 & Identifica endpoints relevantes para herramienta específica \\
\hline
3 & Configura autenticación y permisos según requerimientos de seguridad \\
\hline
4 & Implementa cliente que consume endpoints de acceso a datos \\
\hline
5 & Establece streaming en tiempo real via WebSocket según necesidades \\
\hline
6 & Configura formato de datos compatible con herramienta de destino \\
\hline
7 & Implementa pipeline de transformación de datos si necesario \\
\hline
8 & Configura manejo de errores y reconexión automática \\
\hline
9 & Establece monitorización de calidad de datos transferidos \\
\hline
10 & Valida integridad de integración con datos de prueba \\
\hline
11 & Implementa logging y métricas para monitorización operacional \\
\hline
12 & Despliega integración con documentación de mantenimiento \\
\hline
\end{tabular}
}
\caption{Flujo principal del caso de uso CU-06 - Tercera iteración}
\label{tab:cu06_it3_flujo_principal}
\end{table}

\subsection{Flujos alternativos}
\label{subsec:cu06_it3_flujos_alternativos}

\begin{table}[H]
\centering
\resizebox{\textwidth}{!}{
\rowcolors{2}{white}{udcgray!25}
\begin{tabular}{|p{3cm}|p{11cm}|}
\hline
\rowcolor{udcpink!25}
\textbf{Flujo} & \textbf{Descripción} \\
\hline
\textbf{FA-01: Integración MNE-Python} & 
\begin{minipage}[t]{10.5cm}
\begin{itemize}
    \item En paso 6, configura conversión directa a formatos MNE
    \item Utiliza servicios MNE del sistema para conversión en tiempo real
    \item Mantiene metadatos de electrodos y frecuencia de muestreo
    \item Facilita análisis inmediato en pipeline MNE existente
\end{itemize}
\end{minipage} \\
\hline
\textbf{FA-02: Integración EEGLAB} & 
\begin{minipage}[t]{10.5cm}
\begin{itemize}
    \item Configura exportación en formato .set compatible con EEGLAB
    \item Incluye metadatos de canales en estructura apropiada
    \item Mantiene información de eventos como estructura EEGLAB
    \item Proporciona scripts MATLAB para importación automatizada
\end{itemize}
\end{minipage} \\
\hline
\textbf{FA-03: Integración LabStreamingLayer} & 
\begin{minipage}[t]{10.5cm}
\begin{itemize}
    \item Implementa outlets LSL desde streams internos del sistema
    \item Configura metadatos XML según especificaciones LSL
    \item Mantiene sincronización temporal precisa con reloj LSL
    \item Facilita integración con ecosistema LSL existente
\end{itemize}
\end{minipage} \\
\hline
\end{tabular}
}
\caption{Flujos alternativos del caso de uso CU-06 - Tercera iteración}
\label{tab:cu06_it3_flujos_alternativos}
\end{table}

\subsection{Flujos de excepción}
\label{subsec:cu06_it3_flujos_excepcion}

\begin{table}[H]
\centering
\resizebox{\textwidth}{!}{
\rowcolors{2}{white}{udcgray!25}
\begin{tabular}{|p{3cm}|p{11cm}|}
\hline
\rowcolor{udcpink!25}
\textbf{Excepción} & \textbf{Manejo} \\
\hline
\textbf{FE-01: API rate limiting} & 
\begin{minipage}[t]{10.5cm}
\begin{itemize}
    \item Cliente supera límites de frecuencia de peticiones
    \item Sistema implementa backoff exponencial automático
    \item Proporciona información sobre límites y tiempo de espera
    \item Permite configurar prioridad de endpoints según criticidad
\end{itemize}
\end{minipage} \\
\hline
\textbf{FE-02: Transformación datos fallida} & 
\begin{minipage}[t]{10.5cm}
\begin{itemize}
    \item Error en pipeline de transformación de formato
    \item Sistema identifica paso específico que causa fallo
    \item Preserva datos originales antes de transformación
    \item Ofrece transformación alternativa o formato raw como fallback
\end{itemize}
\end{minipage} \\
\hline
\textbf{FE-03: Pérdida sincronización temporal} & 
\begin{minipage}[t]{10.5cm}
\begin{itemize}
    \item Deriva temporal detectada entre sistema y herramienta externa
    \item Sistema recalibra sincronización automáticamente
    \item Documenta deriva y correcciones aplicadas
    \item Alerta sobre período de datos potencialmente afectados
\end{itemize}
\end{minipage} \\
\hline
\textbf{FE-04: Incompatibilidad versión herramienta} & 
\begin{minipage}[t]{10.5cm}
\begin{itemize}
    \item Herramienta externa actualizada a versión incompatible
    \item Sistema detecta cambios en formato o protocolo esperado
    \item Proporciona adaptador legacy para versiones anteriores
    \item Sugiere actualización de integración a nueva versión
\end{itemize}
\end{minipage} \\
\hline
\end{tabular}
}
\caption{Flujos de excepción del caso de uso CU-06 - Tercera iteración}
\label{tab:cu06_it3_flujos_excepcion}
\end{table}